\chapter{Lecture 2: New Revolutions in Particle Physics: Standard Model}

These notes follow Leonard Susskind's second lecture in the supplementary Standard Model sequence, with curation by LazyingArt LLC. The lecture begins with a warning and a promise. Quantum chromodynamics is far too large a subject for a single meeting, so we will not attempt its deep formal mathematics here. What we will do instead is extract the basic picture. The path is deliberate: QCD is first framed by analogy with QED, then the lecture backs up into the mathematics of spin, and only after that uses the same machinery to introduce isospin, the \(\Delta\) puzzle, color, gluons, and the first qualitative form of QCD.

\section{QCD by Analogy, and Why We Must Revisit Spin}

Susskind opens by placing quantum chromodynamics beside quantum electrodynamics. QED is the theory of electrons and photons; in the atomic setting it is also the theory of the Coulomb force, the field that binds electrons to nuclei. In the same sense, QCD is the theory of quarks and gluons, and of the composite things that quarks and gluons make. Those composites are hadrons. Protons and neutrons are baryons, with net baryon number; mesons are quark--antiquark pairs; and in either case the gluons are part of the binding medium. The lecture wants us to see that QCD is not merely a list of particles. It is a theory of constituents and of the force that holds them together.

But before QCD proper begins, the lecture steps backward. Why? Because the mathematics of spin will come back twice: first almost directly as isospin, and then again, in a looser but still structurally related way, as color. The lecture explicitly refuses a group-theory treatment. We are told that group theory is really lurking underneath, but the working language will be angular momentum, because that language is already familiar and already rich enough for the combinatorics we need.

\section{Angular Momentum Multiplets and Two Spin-Half Particles}

The reminder begins with the commutator algebra of angular momentum:
\begin{equation}
[L_x,L_y]=i\hbar L_z,
\end{equation}
together with the cyclic companions
\begin{equation}
[L_y,L_z]=i\hbar L_x,
\qquad
[L_z,L_x]=i\hbar L_y.
\end{equation}
The lecture immediately emphasizes what this means operationally. The physics is rotationally symmetric, but our description does not have to display that symmetry at every step. We choose a preferred axis, conventionally the \(z\)-axis, and we ask for the component of angular momentum along that axis. This is not symmetry breaking in nature; it is only the expedient choice that makes the representation manageable. Since the components do not commute, we cannot measure them all at once. So we focus on one of them.

From this point of view a spin multiplet is organized by the largest value attained by the \(z\)-component. The lecture says \(L\) aloud, while the blackboard uses \(\ell\); to avoid confusion with the operators \(L_x,L_y,L_z\), we will write the multiplet label as \(\ell\). The familiar counting rule is then
\begin{equation}
m_z=\ell,\ell-1,\ldots,-\ell,
\qquad
\#\text{ of independent states}=2\ell+1.
\end{equation}
That is the key fact the lecture wants us to carry forward.

Now let us put together two spin-half particles, ignoring orbital angular momentum entirely. The product basis is
\begin{equation}
|\uparrow\uparrow\rangle,\qquad
|\uparrow\downarrow\rangle,\qquad
|\downarrow\uparrow\rangle,\qquad
|\downarrow\downarrow\rangle.
\end{equation}
The maximum \(m_z\) is \(+1\), attained when both spins point up, and the minimum is \(-1\), attained when both point down. But there is already something to notice: there are two distinct basis states with \(m_z=0\), namely \( |\uparrow\downarrow\rangle \) and \( |\downarrow\uparrow\rangle \). This is precisely where the lecture stops being a review and becomes a problem.

\begin{figure}[h]
\centering
\includegraphics[width=0.78\textwidth]{lecture_02_figure_02.png}
\caption{Symmetric mixed-spin state highlighted on the board. The frame preserves the moment at which the mixed-spin sum is singled out before the normalization is written explicitly.}
\end{figure}

The board frame is useful not only because the mixed-spin sum has been circled, but also because the earlier spin algebra still remains visible at the top of the board. One can literally see the lecture passing from the commutator and the \(2\ell+1\) counting rule to the first nontrivial example. In cleaned notation the circled state is
\begin{equation}
\frac{|\uparrow\downarrow\rangle+|\downarrow\uparrow\rangle}{\sqrt2}.
\end{equation}
The factor \(1/\sqrt2\) is the standard normalized completion. It is not yet visible in the earlier frame, but the transcript makes clear that this is the next move in the lecture.

\subsection{Question \& Answer}

\paragraph{Question.} If two spin-half particles give four basis states, why do they not form a single spin-one multiplet?

\paragraph{Answer.} Because a spin-one multiplet has only three states:
\[
m_z=+1,\quad 0,\quad -1.
\]
The two aligned states already account for the \(m_z=\pm1\) members. So a spin-one multiplet can contain only one linear combination from the two-dimensional \(m_z=0\) sector. There is still one orthogonal combination left over. It cannot be spin two, because that would require five states altogether. It cannot be anything half-integral, because the total spin of two spin-half particles is integral. The only remaining possibility is spin zero.

\section{Triplet, Singlet, and the Meaning of the Symmetric State}

The lecture now resolves the counting puzzle exactly in that way. We diagonalize the \(m_z=0\) subspace into symmetric and antisymmetric combinations, and the four product states reorganize themselves as
\begin{align}
\frac12\otimes\frac12 &= 1\oplus 0, \\
\ell=1:&\qquad
|\uparrow\uparrow\rangle,\qquad
\frac{|\uparrow\downarrow\rangle+|\downarrow\uparrow\rangle}{\sqrt2},\qquad
|\downarrow\downarrow\rangle, \\
\ell=0:&\qquad
\frac{|\uparrow\downarrow\rangle-|\downarrow\uparrow\rangle}{\sqrt2}.
\end{align}
This is a familiar identity, but in the lecture it is not merely quoted. It is forced by counting, and then interpreted physically.

\begin{figure}[h]
\centering
\includegraphics[width=0.80\textwidth]{lecture_02_figure_03.png}
\caption{Triplet and singlet decomposition of two spin-\(\tfrac12\) states. The board layout itself matters: the three \(\ell=1\) states are grouped together, while the antisymmetric \(\ell=0\) state is set apart below.}
\end{figure}

The essential interpretive point is the one Susskind presses hardest: the symmetric mixed-spin state has \(m_z=0\), but it is not spin zero. It belongs to the \(\ell=1\) triplet. That is why it rotates into the \(m_z=\pm1\) states when the axis is changed. Zero along the chosen axis is not the same thing as zero total angular momentum. Only the antisymmetric combination has total angular momentum zero:
\begin{equation}
\frac{|\uparrow\downarrow\rangle-|\downarrow\uparrow\rangle}{\sqrt2}
\quad \longleftrightarrow \quad
\ell=0.
\end{equation}

The lecture then sharpens the distinction with a physical test. Suppose we ignore every conservation law except angular momentum. Which of these states could simply disappear? The aligned states clearly cannot. The symmetric \(m_z=0\) state also cannot, because although its \(z\)-component vanishes, it still belongs to the spin-one multiplet. The only combination that could disappear on angular-momentum grounds alone is the antisymmetric singlet. In that sense the singlet behaves, as Susskind says, like empty space as far as angular momentum is concerned.

This also clarifies a question from the room. If we measure the symmetric mixed-spin state along the \(z\)-axis, we obtain zero. But rotate the axis, and the same state reveals the nonzero angular momentum appropriate to the triplet. The lecture wants us to keep that distinction in mind, because exactly this sort of counting-with-interpretation will reappear when spin is replaced by isospin.

\section{Three Spins as the Template for Isospin}

At this point the lecture pauses and asks, in effect, why we are doing all this. Before answering, it runs the same counting argument once more for three spin-half particles. There are
\begin{equation}
2\times 2\times 2=8
\end{equation}
product states altogether. If all three spins point up, the maximum \(m_z\) is \(+\tfrac32\), so a spin-\(\tfrac32\) multiplet must be present. That multiplet contains
\begin{equation}
2\left(\frac32\right)+1=4
\end{equation}
states. Four states therefore remain. They cannot form another spin-\(\tfrac32\) multiplet, because the \(\pm\tfrac32\) extrema have already been used. So the only possibility is spin-\(\tfrac12\), and since four states remain, spin-\(\tfrac12\) appears twice:
\begin{equation}
\frac12\otimes\frac12\otimes\frac12
=
\frac32\oplus\frac12\oplus\frac12.
\end{equation}
The lecture's phrase is: let us keep that in mind.

Now the reason for the spin review arrives. What Susskind really wants to discuss is isospin, the first important internal analogue of spin in particle physics. It does not refer to the rotation of physical space. It refers to a two-state internal structure, mathematically isomorphic to the spin-half system. The secure point, especially where the transcript becomes noisy, is that isospin is an internal copy of the same mathematics, not a literal spatial spin.

The justification is physical, not merely formal. The natural scale of hadron physics is a few hundred MeV. On that scale the up and down quarks are both very light, and, more importantly, nearly degenerate. Electromagnetic effects also distort the symmetry only slightly in the hadronic setting we are considering. So as a first approximation we may neglect the mass difference between \(u\) and \(d\) and ignore their electric charges. Then the pair
\[
u,\qquad d
\]
behaves as a doublet under an approximate internal symmetry. In that approximation we may treat the up and down quarks as the analogue of spin up and spin down.

Historically isospin predates quarks; the lecture briefly recalls that it goes back to Heisenberg and to the older interchange of proton and neutron. But here it is introduced in the modern order: the light quark pair provides the internal two-state system, and from there the old isospin language becomes natural again.

\section{From Spin to Isospin: Up and Down Quarks, Proton and Neutron}

Once the \(u\)--\(d\) doublet is in place, the next step is almost forced. A single quark is not the simplest object we observe in the laboratory. The simplest hadrons in the discussion are three-quark states, namely the proton and the neutron. By direct analogy with the three-spin problem, three light quarks can make an isospin-\(\tfrac12\) sector and an isospin-\(\tfrac32\) sector. The lecture chooses to begin with the milder case, the isospin-\(\tfrac12\) doublet.

The two states are
\begin{equation}
p\sim uud,
\qquad
n\sim udd.
\end{equation}
This is the hadronic counterpart of the spin-half multiplet. Two states only, related by the interchange \(u\leftrightarrow d\).

Their charges follow from
\begin{equation}
Q(u)=+\frac23,
\qquad
Q(d)=-\frac13.
\end{equation}
Thus
\[
Q(p)=\frac23+\frac23-\frac13=+1,
\qquad
Q(n)=\frac23-\frac13-\frac13=0.
\]
The small mass difference between proton and neutron is one of the lecture's practical signs that the symmetry is only approximate, but still good enough to organize the strong-interaction multiplets.

At this point the lecture briefly becomes more precise than it wants to remain. The quarks are labeled by positions \(1,2,3\), and the board writes an antisymmetrized factor involving the two up quarks:
\begin{equation}
d_1\bigl(u_2u_3-u_3u_2\bigr).
\end{equation}

\begin{figure}[h]
\centering
\includegraphics[width=0.58\textwidth]{lecture_02_figure_04.png}
\caption{Antisymmetrized two-\(u\)-quark factor in the proton state. This is documentary evidence for the board expression only; it is not the full proton wavefunction.}
\end{figure}

The lecture immediately tells us not to lean too hard on this formal step. It is being precise, but only locally. The board shows one antisymmetrized factor, not the full proton wavefunction. What matters is the analogy with the two-spin singlet: an antisymmetric pair behaves like the zero-angular-momentum combination in the earlier discussion, and together with the third quark one arrives at an overall isospin-\(\tfrac12\) state. Susskind then says, in effect, never mind the fancy part; the important point is that the proton and neutron form the isospin-half doublet.

If one wished to normalize the displayed two-\(u\)-quark factor in the usual way, one would insert a factor of \(1/\sqrt2\), but that factor is not written on the board and is not needed here. The chapter should preserve the actual board evidence and no more.

\section{The Delta Multiplet and the Pauli Puzzle}

Now the lecture shifts to the stronger case. If three light quarks can make isospin \(\tfrac12\), they can also make isospin \(\tfrac32\). The analogue of the spin-\(\tfrac32\) quartet is the \(\Delta\) multiplet:
\begin{equation}
\Delta^{++}\sim uuu,\qquad
\Delta^{+}\sim uud,\qquad
\Delta^{0}\sim udd,\qquad
\Delta^{-}\sim ddd.
\end{equation}
All four states are related by the approximate isospin symmetry and therefore have nearly the same mass, up to the small effects of the \(u\)--\(d\) mass difference.

At the same time these states are not just the proton and neutron in disguise. Their ordinary spin is different. The proton and neutron have actual spin \(\tfrac12\); the \(\Delta\) states have spin \(\tfrac32\). In the lecture this means that the three ordinary spins are aligned, and the three isospins are aligned as well. The quartet therefore makes very concrete the combinatorial pattern that we first met in the abstract three-spin problem.

Their charges are
\begin{align}
Q(\Delta^{++})&=+2, &
Q(\Delta^{+})&=+1, &
Q(\Delta^{0})&=0, &
Q(\Delta^{-})&=-1.
\end{align}
The lecture quotes the proton and neutron mass scale as roughly \(940\) MeV and the \(\Delta\) mass scale as roughly \(1200\) MeV. That is enough to make the \(\Delta\) unstable.

A worked example is the decay of the doubly charged member. Start from
\begin{equation}
\Delta^{++}\sim uuu.
\end{equation}
As the quarks pull apart, a quark--antiquark pair may appear. The lecture first toys with a \(u\bar u\) insertion, then rejects it on energetic grounds and chooses \(d\bar d\). That gives
\begin{align}
uuu &\longrightarrow uud + u\bar d \\
&= p + \pi^+.
\end{align}
Charge conservation is immediate:
\[
Q(p)=+1,
\qquad
Q(\pi^+)=Q(u)+Q(\bar d)=\frac23+\frac13=+1.
\]
So the final state again has total charge \(+2\). The mass balance is also consistent with the lecture's rough numbers:
\[
1200\ \text{MeV} > 940\ \text{MeV}+140\ \text{MeV},
\]
so the decay has phase space available. This is why the \(\Delta\) states are real but very short-lived resonances.

\subsection{Question \& Answer}

\paragraph{Question.} How can a \(\Delta^{++}\) made of three up quarks avoid violating Fermi statistics?

\paragraph{Answer.} This is the central tension of the lecture. In \(\Delta^{++}\) the three quarks are all \(u\), and because the baryon has spin \(\tfrac32\), their ordinary spins are aligned as well. So if spin and flavor were the whole story, we would seem to have three identical fermions in the same state. That is forbidden by the Pauli principle. The lecture presents this not as a minor technical issue but as a historical contradiction that had to be resolved. The resolution is that quarks must carry one more label, hidden at first from view, so that the three quarks in the \(\Delta^{++}\) are not identical in every respect. That extra label is color.

The lecture even pauses to note why the \(\Delta\) was so important historically. In the proton and neutron one could imagine the problem being disguised by more complicated symmetrization. In the \(\Delta^{++}\) the contradiction is blunt: three up quarks, parallel spin, parallel isospin, and something is plainly wrong unless a new quantum number is introduced.

\section{Color, Gluons, and the Minimal Picture of QCD}

Color now enters as the hidden label demanded by the \(\Delta\) puzzle. The lecture uses the conventional names
\[
r,\qquad g,\qquad b,
\]
with corresponding anti-colors
\[
\bar r,\qquad \bar g,\qquad \bar b.
\]
The names are arbitrary. They do not refer to visible color, and the three colored versions of a given quark have the same mass and the same obvious properties. But they are distinct states, and that distinction is what rescues Fermi statistics.

Here the lecture makes a historical analogy that is worth preserving. Spin itself had once played this role in atomic physics: when two electrons seemed to occupy the same orbital state, the new quantum number of spin resolved the Pauli tension. Color plays the same conceptual role a second time. The strategy is the same: attach another label so that the particles are not really identical in all degrees of freedom. In the \(\Delta^{++}\), one red, one green, and one blue quark are enough to remove the contradiction.

The lecture then pivots from the existence of color to the dynamics built from it. Before doing so, it reminds us how economical QED is. The whole theory can be summarized by one basic vertex: an electron emits or absorbs a photon. For bookkeeping purposes one may even pretend that a photon behaves in some respects like an electron--positron pair. Not because it is literally made of one, but because it shares the neutral charge bookkeeping and the spin-one structure. That fiction is not especially useful in QED, but it becomes useful in QCD.

A quark is therefore organized as a three-component color state,
\begin{equation}
q=
\begin{pmatrix}
r\\
g\\
b
\end{pmatrix},
\qquad
\bar q=
\begin{pmatrix}
\bar r & \bar g & \bar b
\end{pmatrix}.
\end{equation}
A gluon is then introduced, at the lecture's level, as a color--anticolor bookkeeping object:
\begin{equation}
q\,\bar q \sim
\begin{pmatrix}
r\bar r & r\bar g & r\bar b\\
g\bar r & g\bar g & g\bar b\\
b\bar r & b\bar g & b\bar b
\end{pmatrix}.
\end{equation}
This suggests nine gluons at first sight. The lecture then immediately adds the caveat that one diagonal superposition is not an independent state. The full group-theoretic derivation is postponed, but the minimal cleaned statement is
\begin{equation}
3\otimes \bar 3 = 8\oplus 1.
\end{equation}
The lecture's phrasing is more heuristic: it looks like nine, but really one particular diagonal sum is ``a nothing.'' That is enough for the present level.

The primitive quark--gluon interaction is the color-changing emission process
\begin{equation}
q^a \longrightarrow q^b + g^{a\bar b}.
\end{equation}
For example,
\begin{equation}
r \longrightarrow g + (r\bar g).
\end{equation}
This is the QCD analogue of photon emission in QED, but now the emitted quantum itself carries color bookkeeping labels. That is why the lecture keeps emphasizing the rule: follow the lines. Track the color through the diagram. If a line turns around in time, reinterpret it as an antiparticle. At this level that is essentially the whole computational grammar.

\begin{figure}[h]
\centering
\begin{tikzpicture}[scale=1.0]
\draw[->] (0,0)--(2,0) node[midway,above] {$r$};
\draw[->] (2,0)--(4,1) node[midway,above] {$g$};
\draw[dashed] (2,0)--(4,-1) node[midway,below] {$r\bar g$};
\fill (2,0) circle (1.5pt);
\end{tikzpicture}

\bigskip

\begin{tikzpicture}[scale=1.0]
\draw[dashed] (0,0)--(2,0) node[midway,above] {$r\bar b$};
\draw[dashed] (2,0)--(4,1) node[midway,above] {$r\bar g$};
\draw[dashed] (2,0)--(4,-1) node[midway,below] {$g\bar b$};
\fill (2,0) circle (1.5pt);
\end{tikzpicture}
\caption{Schematic QCD vertices in the spirit of the lecture. Above: a quark changes color and emits a gluon. Below: a gluon turns into two gluons. The labels are bookkeeping labels, and the operative rule is to follow the color lines.}
\end{figure}

The upper diagram is the direct analogue of the QED vertex. The lower diagram is the new ingredient. Because photons do not carry electric charge, they do not interact with one another at the primitive level. Gluons do. The lecture constructs this by treating the gluon as a quark--antiquark-like color flow and allowing that flow itself to split. The result is a three-gluon vertex. Once that exists, gluons can be exchanged not only between quarks but also between gluons.

That is the decisive novelty of QCD. Quarks bind by exchanging gluons, but gluons themselves can also bind through gluon exchange. The lecture emphasizes three consequences. First, QCD can contain objects with no quarks at all, only gluonic degrees of freedom. Second, the dynamics of the gluon field is non-linear: one gluon wave does not simply pass through another without effect, and even different parts of a single gluon field configuration can act back on one another. Third, the bookkeeping is more intricate than in electrodynamics, which is why the lecture does not try to canonize every color-labeled sketch on the board. The stable lesson is the line-following rule, not every transient chalk label.

The lecture closes with a heuristic coda on running parameters. Imagine a small mass \(m\) attached to a soft extended blob. Push very slowly, and the whole structure moves together; the effective inertia is the mass of the full composite. Shake very rapidly, and the extended part does not have time to respond; then one probes only the small core. The observed mass depends on the scale of the probe. This is the lecture's route to the statement that the effective parameters of a quark are scale dependent:
\begin{equation}
m_q=m_q(\mu),
\qquad
Q_q=Q_q(\mu).
\end{equation}
This is not a renormalization-group derivation. It is a physical analogy, meant only to explain why the ``mass of a quark'' and even the effective charge one attributes to it depend on how hard one probes the hadron.

The lecture deliberately refuses to close the subject here. It tells us instead what comes next: confinement, hadron structure, and the deeper reason color is never seen in isolation. For the moment the lesson is more modest. QCD works by color, gluons are the carriers of that interaction, and unlike photons they interact with one another.

\section{Summary}

The chapter has followed the lecture's intended motion from analogy to contradiction to resolution. QCD was first introduced by comparison with QED, but the real conceptual preparation was the return to spin. From the commutators and the \(2\ell+1\) rule we were led to the two-spin puzzle, then to the triplet and singlet, then to the three-spin decomposition. Only after that did the lecture reinterpret the same mathematics as isospin.

With isospin in hand, the proton and neutron emerged as the isospin-\(\tfrac12\) doublet, while the \(\Delta\) quartet supplied the isospin-\(\tfrac32\) sector and, more importantly, the Pauli puzzle that forces color into the theory. Once color is admitted, the lecture can finally pass to QCD proper: color-labeled quarks, color--anticolor gluon states, the basic quark--gluon vertex, the three-gluon interaction, and the non-linear dynamics that distinguishes QCD sharply from electrodynamics.

The chapter therefore ends exactly where the lecture wants to end: with a mathematically guided intuitive picture, not with a completed theory of hadrons. The next step is confinement. Here the purpose was to see why color had to be invented, what kind of force carrier the gluon is, and how the first narrative of QCD is assembled.